% s-114-p114-m003; original PDF 114, 115
\Aphorism{s-114-p114-m003}{130}
Werden wir die Vorurtheile des Augenscheins los. Was ist das Individuum, wenn nicht eine Summe von Vorstellungen anderer Individuen? Wir dienen einander zum Verzehren; der Geschmack auf der Zunge dessen, der uns verzehrt, das sind „wir“. Und den geistigsten Geschmack, das vollendete Raffinement im Ausdichten des Anderen hat die Liebe, weil sie auf der unendlich reichen Instinctwelt beruht und aus unerschöpflicher Tiefe schöpft. Wir sind am meisten wir selbst, am feinsten und vielfachsten differentiirt in der Phantasie dessen, der uns liebt; in ihr füllen wir uns mit tausend Inhalten — während die Gleichgültigkeit der Anderen uns nur auf irgend etwas Grobes, Dumpfes, Allgemeines hin prüft. Für die Menge sind wir einfach salzig, bitter oder süss, für Jemanden, der uns liebt, werden wir zum polyphonen Geschmacksaccord, zum Inbegriff von tausend Geschmacksnüancen, zur delicaten Besonderheit eines \Foreign{Ego compositum}.
\par

